\section{Calculus and background results}
\label{sec:preliminaries}

We collect the analytic conventions and the precise approximation
facts used below. Throughout this section \((M,d,m)\) is an
\(\RCD(K,N)\) space with \(N<\infty\), full support, and the
usual locally finite reference measure. All the analytic results
recalled here also hold under \(\RCD^*(K,N)\). We use them under the
unreduced hypothesis of Theorem~\ref{main:equivalence}.

\subsection{Test calculus and localization}

Besides \(\Test\), defined in \eqref{eq:test-class}, it is useful to
write
\[
\Test^\infty(M)=\{g\in\Test(M):\Delta g\in L^\infty(M)\}.
\]
The superscript refers to the additional Laplacian bound; it does
not mean classical smoothness. A function is a local test function
on an open set if it agrees on every relatively compact smaller
open set with a global test function. We write \(H_g\) for
\(\Hess g\), and identify the covector \(H_g(V,\cdot)\) with a
vector when using the notation \(H_gV\).

The following facts belong to the standard second-order
\(\RCD\) calculus \cite{Gigli2018,Gigli2019,Savare2014}.
Both test classes are algebras. A test function has an \(L^2\)
Hessian, and for test fields
\[
X\in L^2(TM)\cap L^\infty(TM),\qquad
\nabla X\in L^2(T^*M\otimes TM),\qquad
\divg X\in L^2(m).
\]
The connection is metric-compatible and torsion-free. In particular,
inner products of bounded test fields belong to \(W^{1,2}\),
and their differentials obey the usual product rule.
For gradients,
\[
\nabla_{\nabla g}\nabla g=\nabla\bigl(|\nabla g|^2/2\bigr).
\]
The curvature functional \eqref{eq:curvature-definition} has the
algebraic symmetries of Riemannian curvature and is tensorial over
\(\Test\), including
\begin{align}
\label{eq:curvature-symmetries}
R(X,Y,Z,W)&=-R(Y,X,Z,W)=R(Z,W,X,Y),\\
\label{eq:curvature-tensoriality}
R(aX,Y,Z,W)(f)&=R(X,Y,Z,W)(af)
\end{align}
for test functions \(a,f\), with the corresponding identity in every
slot \cite[Proposition~2.7]{Gigli2019}.

\begin{lemma}[Cutoffs and approximation]
\label{pre:approximation}
The following approximations are available.
\begin{enumerate}
\item For a compact set \(E\) in an open set \(U\), there is
\(\chi\in\Test^\infty\) with \(0\le\chi\le1\), equal to one on a
neighborhood of \(E\), and compactly supported in \(U\).
There are also exhausting cutoffs \(\chi_R\in\Test^\infty\),
equal to one on \(B_R(o)\), supported in \(B_{R+1}(o)\), whose
gradient and Laplacian bounds are independent of \(R\).
\item For \(g\in D(\Delta)\),
\begin{equation}
\label{eq:pre-hessian-estimate}
\|H_g\|_2^2\le
\|\Delta g\|_2^2+\max\{-K,0\}\|\nabla g\|_2^2.
\end{equation}
Compactly supported functions in \(\Test^\infty\) form a core
for \(\Delta\).
\item On a fixed compact region, a function \(g\in\Test\) can be
approximated by compactly supported \(g_j\in\Test^\infty\) with
uniformly bounded gradients, with
\[
g_j\to g,\qquad \nabla g_j\to\nabla g,\qquad
\Delta g_j\to\Delta g,\qquad H_{g_j}\to H_g
\quad\hbox{strongly in }L^2
\]
on that region.
\item A compactly supported Lipschitz function \(\xi\) can be
approximated by compactly supported test functions \(\xi_j\)
with a common compact support and a common Lipschitz bound,
uniformly and strongly in \(W^{1,2}\).
If \(\xi\ge0\), the approximants may be chosen nonnegative.
\end{enumerate}
\end{lemma}

\begin{proof}
The good-cutoff construction and the integrated Bochner estimate give
\textup{(i)} and \eqref{eq:pre-hessian-estimate}; see
\cite[Appendix~A, Lemma~A.2 and Equation~(A.13)]{GigliTamanini2021}.
For the local version of the cutoff construction one uses finitely
many small good cutoffs inside \(U\), and a smooth truncation of their
sum. This keeps the Laplacian bounded and Sobolev by the test
algebra and chain rules.

Let \(P_t\) denote the heat semigroup. For \(h\in L^2\cap L^\infty\),
positive-time regularization gives
\(P_th\in\Test^\infty\)
\cite[Appendix~A, Equation~(A.7)]{GigliTamanini2021}.
Spectral calculus and density of bounded \(L^2\) data show that these
regularizations form a core. For each such regularization \(h\),
\[
\Delta(\chi_Rh)
=\chi_R\Delta h+h\Delta\chi_R
+2\langle\nabla\chi_R,\nabla h\rangle.
\]
The uniform cutoff bounds and the \(L^2\) tails imply convergence
to \(\Delta h\) in \(L^2\), as well as convergence of the functions
and gradients. This proves the compact-core assertion.

For \textup{(iii)}, choose nested cutoffs \(\chi,\zeta\) such that
\(\chi=1\) near the region in question and \(\zeta=1\) near
\(\supp\chi\). Put
\[
g_t=\chi P_t(\zeta g).
\]
The function \(\zeta g\) lies in \(\Test\). The product rule,
spectral continuity and the fixed cutoff bounds give convergence
of \(g_t\) to \(\chi\zeta g\) in the Laplacian graph norm.
The Bakry gradient estimate gives
\[
\|\nabla g_t\|_\infty
\le \|\chi\|_\infty e^{-Kt}\|\nabla(\zeta g)\|_\infty
+\|\nabla\chi\|_\infty\|\zeta g\|_\infty.
\]
Apply \eqref{eq:pre-hessian-estimate} to
\(g_t-\chi\zeta g\) to obtain strong Hessian convergence.
Since \(\chi\zeta g=g\) near the region, \textup{(iii)} follows.

For \textup{(iv)}, use \(\xi_t=\chi P_t\xi\), with an outer
cutoff \(\chi=1\) near \(\supp\xi\).
The heat semigroup converges strongly in \(W^{1,2}\);
the uniform Lipschitz estimate and the small-time heat-kernel
moment bound give uniform convergence on the common compact
support. The cutoff product gives the same conclusions for
\(\xi_t\), and positivity is preserved.
\end{proof}

\begin{remark}
\label{pre:bare-gradients}
A compactly supported test gradient is a member of the global
\(\TestV\): multiply it by a test cutoff equal to one near its
support. This observation is necessary when \(m(M)=\infty\),
because the constant function one then need not lie in \(\Test\).
All uses of a bare test gradient below are local in this sense.
\end{remark}

\begin{lemma}[Lipschitz weights and continuity of the tensor]
\label{pre:tensor-continuity}
For fixed test fields the formula
\eqref{eq:curvature-definition} defines a functional on compactly
supported Lipschitz weights. It has the form
\[
f\longmapsto\int fA+\langle\nabla f,B\rangle\dd m,
\qquad A\in L^1_{\rm loc}(m),\quad B\in L^2_{\rm loc}(TM).
\]
It is continuous under the weight approximation in
Lemma~\ref{pre:approximation}\textup{(iv)}.
More generally, for weights with a common compact support,
uniform convergence of the weights and weak \(L^2\) convergence
of their gradients suffice.
The tensor identities remain valid for these weights.
Moreover, on a fixed compact set the functional is continuous
under uniformly bounded field approximations for which the fields,
their divergences and their covariant derivatives converge
strongly in \(L^2\), with the fields converging almost everywhere
along a subsequence.
\end{lemma}

\begin{proof}
Expand \(\nabla_X(fW)=X(f)W+f\nabla_XW\).
Every coefficient of \(f\) is a product of two \(L^2\) factors
and bounded factors. The coefficients of \(\nabla f\) lie in
\(L^2\). Uniform convergence of \(f_j\) pairs with the fixed
\(L^1\) coefficient, and weak \(L^2\) convergence of their
gradients pairs with the fixed \(L^2\) coefficient.
This proves the weight-continuity assertions.
Approximation then extends the identities.

For the field assertion,
\[
\nabla_{X_j}Y_j-\nabla_XY
=(\nabla Y_j-\nabla Y)(X_j)+\nabla Y(X_j-X).
\]
The first term tends to zero in \(L^2\) by uniform boundedness
of \(X_j\); the second does so by dominated convergence.
The brackets therefore converge in \(L^2\) as well.
Products of the resulting strongly convergent \(L^2\) quantities
converge in \(L^1\). Any remaining uniformly bounded field factor
is handled by almost-everywhere convergence and domination.
This includes the last term of
\eqref{eq:curvature-definition}.
\end{proof}

\subsection{Weak Hessians and Bochner measures}

For \(F\in W^{1,2}_{\mathrm{loc}}(M)\), a compact test gradient \(Y=\nabla g\),
and a compactly supported Lipschitz weight \(\xi\), define
\begin{equation}
\label{eq:weak-hessian}
H_F(Y,Y)(\xi)
=-\int_M\divg(\xi Y)\,\langle\nabla F,Y\rangle\dd m
-\int_M\xi\,\langle\nabla F,\nabla_YY\rangle\dd m.
\end{equation}
Polarization defines the mixed evaluation.
The two terms are integrable: \(\nabla F\), \(\divg(\xi Y)\)
and \(\nabla_YY\) are locally \(L^2\), and \(Y,\xi\) are bounded.
For fixed \(Y,\xi\), the functional is continuous under weak
local \(W^{1,2}\) convergence of \(F\).
When \(F\in D(\Delta)\), this agrees with the usual \(L^2\)
Hessian. A statement \(H_F(Y,Y)\le A\,m\) means the
corresponding inequality on every nonnegative compact Lipschitz
weight. It does not presuppose that \(H_F\) is an \(L^2\) tensor.

For a local test function \(g\), the measure-valued Bochner
inequalities are
\begin{align}
\label{eq:bochner-hessian}
\Gamma_2(g)
&:=\Delta\bigl(|\nabla g|^2/2\bigr)
-\langle\nabla g,\nabla\Delta g\rangle m
\ge \bigl(K|\nabla g|^2+|H_g|^2\bigr)m,\\
\label{eq:bochner-dimension}
\Gamma_2(g)&\ge
\bigl(K|\nabla g|^2+(\Delta g)^2/N\bigr)m.
\end{align}
Here the first Laplacian is a measure-valued Laplacian.
We use the two inequalities separately or after taking a convex
combination; no inequality with the sum of both squared terms
and coefficient one is asserted.
Our capacity normalization is
\[
\Cap_2(E)=
\inf_{\substack{U\supset E\\U\ \mathrm{open}}}
\ \inf_{\substack{v\in W^{1,2}(M)\\v\ge1\ m\text{-a.e.\ on }U}}
\int_M\bigl(v^2+|\nabla v|^2\bigr)\dd m.
\]
A polar set has zero \(2\)-capacity; \emph{quasi-everywhere} means
outside a polar set.
The Bochner measures do not charge sets of zero \(2\)-capacity.
Their Sobolev multipliers are evaluated using quasi-continuous
representatives \cite[Sections~2.3 and 3.2]{Savare2014}.

More explicitly, if bounded Sobolev multipliers converge strongly
in \(W^{1,2}\), with a common bound and compact support, a
subsequence converges outside a polar set. Dominated convergence
therefore applies to their integrals against each fixed local
Bochner measure. This extends the integration-by-parts identities
from test multipliers to bounded Sobolev multipliers.
The same argument shows that an almost-everywhere lower bound for
a Sobolev multiplier holds for its quasi-continuous representative
outside a polar set. These facts will be used when multiplying
singular Bochner measures in Section~\ref{sec:heat-index}.

The inequalities \eqref{eq:bochner-hessian}--\eqref{eq:bochner-dimension}
also hold for our original test class, without a bounded
Laplacian. One can see the extension locally by applying heat
regularization to a test function. The gradients remain uniformly
bounded, the Hessians converge strongly in \(L^2\), and
\(\nabla\Delta P_tg\to\nabla\Delta g\) in \(L^2\).
The local masses of the positive Bochner remainders are bounded
by testing with a fixed outer cutoff. Passage to the distributional
limit gives the asserted measure inequalities.

\subsection{Geometry, transport and along-curve calculus}

We use the following established facts, stating the restrictions
that are relevant to the proof.

\begin{enumerate}
\item A finite-dimensional full-support \(\RCD\) space is proper
and geodesic, is locally doubling, and satisfies a local
\(L^2\) Poincar\'e inequality and the Sobolev-to-Lipschitz property.
Bounded sets have finite measure. When the dimension parameter is
\(N=n\) and \(m=\mathcal H^n\), almost every point is \(n\)-regular and the
tangent module has dimension \(n\) almost everywhere. For every
bounded region and every fixed sufficiently small upper radius,
\begin{equation}
\label{eq:local-volume}
c\,r^n\le m(B_r(x))\le C\,r^n
\end{equation}
with positive constants uniform for centers in that region.
The upper bound is the noncollapsed density bound; the lower bound
follows from Bishop--Gromov and full support
\cite[Theorem~1.12 and Corollary~2.14]{DPG2018}.

\item Optimal quadratic transport between probability measures
\(\mu,\nu\in\mathcal P_2(M)\), with \(\mu\ll m\), is induced by
a unique dynamical map.
For a fixed real-valued \(c\)-concave potential, \(c=d^2/2\), the geodesic
from almost every point into its \(c\)-superdifferential is unique
\cite[Theorems~1.1 and 1.3]{GigliRajalaSturm2016}.
Between bounded densities with bounded support, the interpolating
densities are uniformly bounded and obey the pointwise
curvature-dimension distortion inequality.
We use the density bound also on subintervals with a singular
endpoint, whenever a separate positive-time compression estimate
has been proved. We do not assume bounded compression at a Dirac
endpoint.

\item A bounded-speed plan with bounded compression is a
Sobolev test plan. For a Sobolev function \(g\), its time
derivative is \(\langle\nabla g,v\rangle\), where \(v\) is the
plan velocity \cite[Theorem/Definition~2.33]{GigliDifferentialStructure}.
For a \(W_2\)-continuous curve with bounded compression, the
Sobolev continuity equation characterizes \(2\)-absolute
continuity and its action, and such curves have a superposition
by absolutely continuous metric paths
\cite[Theorem~3.5 and Proposition~2.1]{GigliHan2015}.
The required \(W_2\)-continuity is verified before using this
criterion in Section~\ref{sec:heat-index}.

\item For a bounded-support optimal geodesic test plan,
Gigli--Tamanini's second-order differentiation formula applies to
\(H^{2,2}\) functions and to vector-field pairings with
\(H_C^{1,2}(TM)\), the closure of \(\TestV(M)\) in the covariant
Sobolev norm \cite[Theorem~5.13]{GigliTamanini2021}.
Every vector field to which we apply this theorem is a compactly
supported test field.
It also applies after bounded measurable restriction of the
plan. This permits passage from integrated identities to
almost-sure time-distributional identities.

\item A compactly supported time-dependent test vector field,
bounded with bounded divergence and integrable Sobolev covariant
derivative, has the regular Lagrangian flow and uniqueness
properties used for the short flows below
\cite[Sections~8 and 9.6]{AmbrosioTrevisan2014}.
The compression estimate is the exponential of the integrated
negative-divergence bound. The corresponding backward flow gives
the inverse estimate. These are Sobolev flows; no spatial
Lipschitz bound on the vector field is assumed.
\end{enumerate}

The more specialized harmonic splitting, tangent-continuity,
heat-kernel and variable-curvature theorems are stated at their
points of use. For the geometric implication from \textup{(A)}
to the metric-measure part of \textup{(B)}, we use the Alexandrov
curvature-dimension and infinitesimal Hilbertianity results
\cite{Petrunin2011,ZhangZhu2010,KMS2001}. The unreduced statement
for all curvature bounds, including Alexandrov spaces with boundary,
is recorded in \cite[Theorem~2.9 and Corollary~2.10]{KapovitchKettererSturm}.
For the converse we use unreduced \(\CD\) as a hypothesis; invoking
analytic results written with \(\RCD^*\) consequently changes
neither \(K\) nor \(N\).
